% !TeX encoding = UTF-8
\documentclass[variant=guia,cover=institutional,mode=screen]{ua-engineering-v6-article}
% !TeX encoding = UTF-8
\UASetUniversity{Universidad Austral}
\UASetFaculty{Facultad de Ingeniería}
\UASetDepartment{Departamento de Computación}
\UASetCampus{Pilar}
\UASetCareer{Ingeniería Informática}
\UASetCommission{Comisión A}
\UASetProfessor{Equipo docente}
\UASetSemester{Segundo cuatrimestre 2026}
\UASetCourse{Sistemas Distribuidos}
\UASetDate{2026}

\UASetClassNumber{07}
\UASetClassTitle{Logs distribuidos y streaming}
\title{Clase 07 --- Guía de ejercicios}
\author{\UAProfessor}
\date{\UADate}
\begin{document}
\maketitle
\section{Propósito de la guía}
Los ejercicios exigen definir la frontera de cada garantía. No se acepta “Kafka da exactly-once” sin especificar qué estado, efecto y sistema participan.
\section{Ejercicios}
\begin{uaexercise}
\textbf{1. Log, cola y base orientada a estado.} Compare las tres abstracciones para un sistema de pedidos. Explique retención, replay, orden y ownership de progreso.
\end{uaexercise}
\begin{uaexercise}
\textbf{2. Key de partición.} Para eventos de pedidos, pagos y cambios de estado, proponga una key que preserve el orden relevante. Analice hot keys y qué orden se pierde.
\end{uaexercise}
\begin{uaexercise}
\textbf{3. Número de particiones.} Un topic recibe 80.000 eventos/s y se estima que cada consumidor procesa 5.000 eventos/s. Proponga una cantidad inicial de particiones y discuta límites del cálculo.
\end{uaexercise}
\begin{uaexercise}
\textbf{4. Consumer group.} Explique qué ocurre si un grupo tiene 12 consumidores y el topic 8 particiones. ¿Qué cambia si se agregan 16 particiones durante operación?
\end{uaexercise}
\begin{uaexercise}
\textbf{5. Position y committed offset.} Construya un ejemplo donde la posición en memoria avance, pero el committed offset quede atrasado. Analice el reinicio.
\end{uaexercise}
\begin{uaexercise}
\textbf{6. At-most-once.} Diseñe un timeline de pérdida por confirmar offset antes de aplicar el efecto. Indique cuándo este trade-off podría ser aceptable.
\end{uaexercise}
\begin{uaexercise}
\textbf{7. At-least-once.} Diseñe un timeline de duplicación por procesar antes de confirmar. Proponga una estrategia de idempotencia para el efecto.
\end{uaexercise}
\begin{uaexercise}
\textbf{8. Consumer lag.} Defina lag y explique por qué el promedio del grupo puede ocultar una partición caliente. Proponga métricas adicionales.
\end{uaexercise}
\begin{uaexercise}
\textbf{9. Rebalance.} Durante una campaña, el autoscaler crea consumidores y desencadena rebalances frecuentes. Explique cómo esto puede empeorar temporalmente el servicio y proponga mitigaciones.
\end{uaexercise}
\begin{uaexercise}
\textbf{10. ACK del producer.} Compare confirmar con el líder frente a esperar réplicas in-sync. Analice latencia, disponibilidad y pérdida ante failover.
\end{uaexercise}
\begin{uaexercise}
\textbf{11. Producer idempotente.} Explique qué duplicados evita y construya un caso de duplicación de negocio que siga siendo posible.
\end{uaexercise}
\begin{uaexercise}
\textbf{12. Exactly-once por frontera.} Defina exactamente una vez para entrega, procesamiento y efecto de negocio. Dé un sistema donde solo una de las tres se cumpla.
\end{uaexercise}
\begin{uaexercise}
\textbf{13. Kafka Streams transaccional.} Describa cómo coordinar input offsets, state store y output topic. Explique qué ocurre al abortar.
\end{uaexercise}
\begin{uaexercise}
\textbf{14. Dual write.} Un servicio persiste una factura y publica InvoiceCreated. Construya las ventanas de falla para ambos órdenes posibles.
\end{uaexercise}
\begin{uaexercise}
\textbf{15. Transactional outbox.} Diseñe tablas, estados y proceso relay para resolver el dual write. Explique por qué el consumidor aún debe ser idempotente.
\end{uaexercise}
\begin{uaexercise}
\textbf{16. Event time.} Un evento ocurre a las 10:00 y llega a las 10:07. Compare el resultado en ventanas por event time y processing time.
\end{uaexercise}
\begin{uaexercise}
\textbf{17. Watermark.} Proponga una tolerancia de eventos tardíos para un tablero de ventas y otra para detección de fraude. Justifique la diferencia.
\end{uaexercise}
\begin{uaexercise}
\textbf{18. Stateful streaming.} Diseñe un join entre pedidos y pagos. Defina key, retención de estado, timeout de join y comportamiento ante evento tardío.
\end{uaexercise}
\begin{uaexercise}
\textbf{19. Caso de campaña.} Diseñe topics, keys, grupos, semánticas y SLOs para inventario, email y analytics durante una promoción masiva.
\end{uaexercise}
\begin{uaexercise}
\textbf{20. Integración D=(P,F,G,S,T).} Aplique el marco a un pipeline de telemetría industrial y proponga una prueba de falla con consumer crash y rebalance.
\end{uaexercise}
\section{Criterios de entrega}
Incluya timelines, keys, offsets y escenarios de falla cuando corresponda. Una solución sólida define orden, progreso, duplicación, efecto y métricas.
\end{document}
